\clearpage
\thispagestyle{empty}
\pdfbookmark[0]{Anzeigen eigener Abhandlungen. Ausgabe und Quellen}{anzeigen-source-notice}
\begin{center}\Large\textbf{Anzeigen eigener Abhandlungen}\end{center}

KI-Übertragung des deutschen Originals in modernes Deutsch mit KI-erstelltem Satz. Grundlage ist Carl Friedrich Gauss, \textit{Werke}, Band V (Göttingen, 1867), mit dem Abteilungstitel und der Anzeige aus den \textit{Göttingischen gelehrten Anzeigen} vom 5. April 1813, Digitalisatblätter 283--291. Die Anzeige bespricht Gauß’ Abhandlung über die Anziehung homogener Ellipsoide.

\noindent\textbf{Übersetzung des französischen Zitats von Lagrange}\par\smallskip
\noindent Ihre Lösungen kann man nur als Meisterwerke der Analysis betrachten, doch darf man sich weiterhin eine direktere und einfachere Lösung wünschen, und die stetigen Fortschritte der Analysis geben Anlass, darauf zu hoffen.

\medskip\noindent\textbf{Verwandte Darstellung}\par\smallskip
\noindent Isaac Todhunter behandelt die lateinische Abhandlung in \textit{A History of the Mathematical Theories of Attraction and the Figure of the Earth}, Bd.~II (1873), §§~1162--1175. Er nennt auch die historische deutsche Anzeige und ihren Nachdruck.

\medskip\noindent\href{https://gdz.sub.uni-goettingen.de/id/PPN236006339}{Deutsches Original (1867)}\quad
\href{https://www.gutenberg.org/ebooks/69832}{Todhunter (1873)}
\clearpage
